\section{Task 2 : Choice of application, diode and frequency}
\label{sec:task2}

\subsection{Sensor and frequency}
For the project, we needed to supply an ultra low power sensor. The choosed component is STTS22H, 
a temperature sensor made by STMicroelectronics. We can justify this choice with its ideals caracteristics 
for energy harvesting. It’s an I2C temperature sensor working between 1.5V and 3.6V and it can be on with 
an ultra low supply current under 120 µA. To harvest ambient energy, we choosed 2.2 GHz frequency. 
This frequency is between 1.5 and 3 GHz wanted in the specifications.

\subsection{CMOS Diodes vs Schottky Diodes}
Harvesting RF signals (which are often very weak, around a microwatt) requires an extremely sensitive and fast AC-DC rectifier.

- CMOS Diodes (P-N Junction): These are based on a junction of P-doped and N-doped semiconductors. 
They have a relatively high threshold voltage (usually between 0.6 V and 0.7 V). 
Additionally, their switching time is limited by the recombination of minority charge carriers. 
Therefore, they are not suitable for rectifying very weak, high-frequency RF signals.

- Schottky Diodes: These use a metal-semiconductor junction. This structure provides a much lower threshold voltage 
(typically between 0.15 V and 0.3 V), allowing the rectification to start even with very low input power. 
Moreover, since the current is carried by majority carriers, there is no delay caused by recombination. 
This enables the ultra-fast switching required at gigahertz (GHz) frequencies.

\subsection{BAV99 and HSMS-28xx series diodes}
The BAV99 is a standard silicon switching diode (P-N). Because of its high threshold voltage 
and large junction capacitance, it would severely weaken the signal at 2.2 GHz. It would also 
require much more input power than an energy harvesting antenna can actually provide.

The HSMS-28xx series consists of Schottky diodes specifically designed for RF detection and energy harvesting. 
They are optimized to offer minimal junction capacitance (Cj) and very low series resistance (Rs), 
which maximizes the RF-to-DC conversion efficiency.

\subsection{Which diode do we need ?}

We compared three models from the HSMS series to find the best diode for our 2.2 GHz target frequency. 
The HSMS-2850 has a low threshold voltage, but it is designed to operate at frequencies below 1.5 GHz. 
Therefore, we cannot use it for our 2.2 GHz application. The HSMS-2820 can work up to 4 GHz, but it has a 
junction capacitance of 0.7 pF. At 2.2 GHz, this capacitance would cause significant insertion losses. 
The HSMS-2860 works up to 5.8 GHz and has a very low junction capacitance of 0.18 pF.

The HSMS-2860 diode is the best choice for our matching and rectifying circuit at 2.2 GHz. 
Its low parasitic parameters will guarantee the best conversion efficiency to power the STTS22H sensor.

\begin{table}[h]
    \centering
    \caption{Schottky diodes compraison for RF energy harvesting applications at 2.2 GHz.}
    \label{tab:diodes_comparison}
    \begin{tabular}{lcccc}
        \toprule
        \textbf{Diode} & \textbf{($I_s$)} & \textbf{($R_s$)} & \textbf{($C_j$)} & \textbf{($F$)}\\
        \midrule
        HSMS-2820 & \SI{22}{\nano\ampere} & \SI{6}{\ohm} & \SI{0.7}{\pico\farad} & up to \SI{4}{\giga\hertz} \\
         HSMS-2850 & \SI{3000}{\nano\ampere} & \SI{25}{\ohm} & \SI{0.18}{\pico\farad} & below \SI{1.5}{\giga\hertz} \\
        HSMS-2860   & \SI{50}{\nano\ampere} & \SI{6}{\ohm} & \SI{0.18}{\pico\farad} & up to \SI{5.8}{\giga\hertz} \\
        \bottomrule
    \end{tabular}
\end{table}